\chapter{Snow products}
\label{cap:snow_products}

The previous chapter introduced the main physical processes governing the snowpack, and showed that snow cover is a highly dynamic system which can be 
investigated by means of several measurements techniques. None of them, however, is able on its own describe every property of the snowpack over a wide 
and morphologically complex region. Manual surveys and automatic weather stations provide accurate local observations, but their spatial 
representativeness is limited, since monitoring networks are often too sparse to capture the variability of the snowpack and are typically 
biased toward low to intermediate elevations. Remote sensing, in contrast, provides the spatial coverage that ground observations lack, but at the 
price of significant trade-offs in resolution, revisit time and the type of variable that can be retrieved. MODIS, for instance, yields daily snow 
cover maps at $\SI{500}{\meter}$ resolution, but provides no hydrological insight, as it only discriminates between snow-covered and 
snow-free pixels. Remote sensing techniques targeting snow water equivalent (SWE), on the other hand, typically operate at resolutions of tens to 
hundreds of kilometers, too coarse to resolve local topography and therefore suitable only for hydrological studies at the continental scale.  

To overcome these limitations and obtain continuous information on snow variables across an entire region, numerical models have been developed to 
simulate the evolution of the snowpack under meteorological forcing. The simplest approaches rely solely on air temperature and total precipitation, regarded 
as the main drivers of snow accumulation and snowmelt. More advanced models solve the mass and energy balance equations (see Chapter \ref{cap:snow_basics}) 
at every grid point, accounting for the effects of complex terrain and for microphysical processes such as snow metamorphism, with the goal of reproducing 
the state of the snowpack as accurately as possible. Their input fields, which must be spatialized beforehand, include incoming solar radiation, air 
temperature, relative humidity and wind speed, among others. 

Some of these models, such as ALPINE3D, are already employed as tools for avalanche forecasting and danger assessment, since the simulations can accurately 
identify weak snow layers that are prone to failure \parencite{ALPINE3D}. Snow products are equally relevant in hydrological applications: an accurate estimate of 
the amount of water stored within the snowpack and of the timing of its release can substantially improve discharge predictions, which in turn are essential 
for the management of hydropower production. Beyond these practical applications, snow models play a fundamental role in climate change research, as they 
allow spatially distributed trend analysis and even make it possible to evaluate the response of the snowpack to different emission scenarios. 

This flexibility, however, comes at a cost. The snow variables resulting from numerical simulations inherit the uncertainties of the meteorological 
forcing, whose fields must themselves be interpolated over complex terrain in order to match the resolution required by the model. Moreover, the modelling 
philosophy behind each product may neglect physical processes whose omission proves detrimental in specific applications: most models are therefore 
task-oriented, and should be used only for the class of problems they were designed to address. Since different models may yield substantially different 
results for the same snowpack, assessing the extent of such disagreement is a prerequisite for producing reliable snow predictions in a given domain. This 
is exactly the aim of this chapter: the following sections introduce and compare the main snow datasets available for the study area of this thesis, 
highlighting their different approaches, the discrepancies that emerge over the investigated region and the inadequacy for trend estimations, thereby 
motivating the need for a new product. 










\section{Study area}

This thesis focuses on the Italian Hydrological portion of the Greater Alpine Region (IH-GAR), comprising all Italian territory located north of $\SI{43}{\degree N}$ 
together with the Swiss portion of the Po basin. The main geological features of this territory are respectively: 
\begin{itemize}[label=$\blacktriangleright$]
    \item \textit{the Italian Alps}, which are the southernmost of the Alpine Arc. They extend for approximately $\SI{1200}{\kilo\meter}$, covering an 
    area of about $\SI{50000}{\kilo\meter\squared}$. The highest peak, Monte Bianco, reaches $\SI{4810}{\meter}$ a.s.l., and summits above 
    $\SI{3000}{\meter}$ a.s.l. are  fairly common throughout the range. Such elevations allow snow to persist throughout summer at the highest altitudes, 
    sustaining 901 glaciers, whose total surface area amounted to roughly $\SI{368}{\kilo\meter\squared}$ in \num{2016}. \parencite{catasto_ghiacciai}.
    \item \textit{the Po Plain}, which is a predominantly flat region shaped by an initial phase of marine sedimentation followed by later fluvial deposition. 
    It is drained by the Po river, Italy's longest, which flows for approximately $\SI{652}{\kilo\meter}$. Together with the Venetian and Friulian plains, 
    it forms the most productive agricultural and industrial area of the country, home to roughly 20 million people.
    \item \textit{the Northern Apennines}, formed by the collision between the western margin of the African Promontory and the Sardinia-Corsica block. 
    Their highest peak reaches only $\SI{2165}{\meter}$ a.s.l., insufficient to prevent the complete melt of the seasonal snowpack
    \parencite{libro_geografia_IHGAR}. 
\end{itemize}
The high elevations of the Italian Alps, combined with low winter temperatures, allow substantial volumes of water to be stored in solid form. Snow also 
accumulates over the Northern Apennines, although in far smaller amounts and for shorter periods, since the peaks are substantially lower. These 
differences in snow cover give rise to distinct hydrological regimes: up to 50\% of the annual streamflow of rivers rising in the Alpine Arc consist 
of meltwater \parencite{snow_drought}, whereas rivers draining the Apennines rely predominantly on rainfall, and often experience low flow, or even run 
completely dry, during summer. The complex orography of the study area, examined in greater detail in Figure \ref{fig:study_area}, also plays a 
significant role in governing the climate of Northern Italy. 
\begin{figure}[htbp]
    \centering
    \includegraphics[width=\linewidth]{snow_products/study_area.png}
    \caption{Topography and elevation distribution of the study area. The frequency histogram is based on 100 
    meters wide elevation bins. Elevation data were obtained from ALOS AW3D30 global digital elevation model (DEM).}
    \label{fig:study_area}
\end{figure}
Indeed, the Italian Alps act as a barrier against cold air masses from the north, while the Apennines shield the Po Valley from the milder influence of 
the Tyrrhenian Sea. Two main climate types can thus be investigated within the study area:
\begin{itemize}[label=$\blacktriangleright$]
    \item \textit{the Alpine climate}, which prevails over the Alps and the Northern Apennines and is characterized by low nighttime and winter 
    temperatures, with precipitation peaking in summer;
    \item \textit{the Po Valley climate}, with cold winters, hot summers, and precipitation concentrated during spring and autumn 
    \parencite{clima_IHGAR}.
\end{itemize}

This complex climatological and topographical setting gives rise to a substantial variability in the snowpack, which makes the IH-GAR region an ideal test 
bed for comparing the performance of snow models and remote sensing products in reproducing the actual state of the snow cover. The strong dependence on 
orography can be investigated by means of satellite observations alone, whereas the accurate streamflow estimates required by agriculture and hydropower 
production call for more complex and physically based modelling. To date, three spatially distributed snow datasets are available over the area: MODIS, 
IT-SNOW and a product specifically developed for the mountainous portion of the Po River District. None of them matches the IH-GAR boundaries exactly, but 
all overlap with the investigated domain to a large extent. Each dataset is discussed in a dedicated section below. The main results presented hereafter 
were computed as a part of the spin-up work of this thesis, which was made possible by the public online availability of the three datasets. In each case, 
the methodological framework outlined in the corresponding literature was adopted.










\section{MODIS}
\label{sec:MODIS}

Owing to its multispectral configuration, MODIS has already proven to be a key tool for mapping snow cover extent and duration across a variety of 
regions, such as the Karakoram \parencite{MODIS_KARAKORAM}. One of the most recent studies based on this sensor is devoted to the Italian territory 
\parencite{MODIS_ITA} and thus covers almost the entirety of the IH-GAR. The resulting dataset, curated by the Department of Enviromental Science and 
Policy of the University of Milan, provides daily snow cover maps at the national scale and, in its latest version, spans the period from \num{24} 
February \num{2000} to \num{30} September \num{2025}, thereby including \num{25} hydrological years. It consists of binary rasters with a spatial 
resolution of approximately $\SI{500}{\meter}$, in which a pixel is assigned a value of 1 when it's at least half-covered in snow, and a value of 0 when 
the ground is bare.



\subsection{Dataset generation and gap-filling strategy}

To obtain binary rasters describing the presence or absence of snow cover, \cite{MODIS_ITA} proposed an operational processing chain 
designed to maximise prediction reliability. The procedure makes use of the MODIS Terra Snow Cover Daily L3 Global $\SI{500}{\meter}$ 
SIN Grid V061 (MOD10A1 V61) and MODIS Aqua Snow Cover Daily L3 Global $\SI{500}{\meter}$ SIN Grid V061 (MYD10A1 V61) products, both 
distributed by the US National Snow and Ice Data Center (NSIDC). In these datasets, an integer value is assigned to each grid cell 
according to the criteria summarised in Table \ref{tab:integer_value}.

\begin{table}[htbp]
  \centering
  \caption{Integer values in NSIDC snow product.}
  \label{tab:integer_value}
  \renewcommand{\arraystretch}{1.3}
  \begin{tabular}{@{} >{\raggedright\arraybackslash}p{2.5cm} >{\raggedright\arraybackslash}p{6.5cm} @{}}
    \toprule
    \textbf{Value} & \textbf{Condition} \\
    \midrule
    0           &  Snow free\\
    1-100            & Snow covered (NDSI value) \\
    200      & Missing data \\
    201          & No decision \\
    211          & Recorded during the nigth \\
    237     & Fresh water \\
    239      & Sea water \\
    250              & Cloud covered \\
    254           & Detector saturated \\
    255               & Fill value \\
    \bottomrule
  \end{tabular}
\end{table}

To address the observation gaps caused by cloud cover or by sensor saturation, \cite{MODIS_ITA} first reduce the complexity of the dataset to three 
surface classes: snow covered, snow free and missing data. Integer values between 1 and 100 are reclassified as snow covered 
whenever the corresponding NDSI exceeded the critical threshold of 0.4, and as snow free otherwise. Pixels aquired at night, 
affected by cloud cover or flagged with an error condition are all labelled as missing data, whereas freshwater pixels are 
reclassified as missing only when freshwater presence has been detected in less than 10 \% of the available observations. Once the 
number of classes has been reduced to three, \cite{MODIS_ITA} propose a gap filling procedure structured into six sequential phases, 
each exploiting different spatial and temporal information to progressively reclassify the missing pixels untill the record is 
complete. These steps, adapted from the previous work of \cite{no_cloud_MODIS} consist of the following:
\begin{itemize}[label=$\blacktriangleright$]
    \item \textit{Terra-Aqua combination}, which exploits the few hour offset between Terra and Aqua overpasses. Since Terra is 
    considered more accurate than its counterpart, the snow map is built primarily from Terra values whenever snow presence or 
    absence is known. If a Terra data is missing, the corresponding Aqua value is examined instead and retained if it is not itself 
    missing. This step assumes that no snowmelt or snowfall occurs between the two overpasses.
    \item \textit{Short range temporal filling}, which relies on combining snow cover information from adjacent dates to try filling 
    the remaining gaps. For each grid point, the values recorded on the previous and following day are considered: if neither is 
    missing and both agree on the pixel classification, the missing value is reclassified accordingly.
    \item \textit{Elevation based filling}, which relies on the vertical extent of snow cover on a given day. For each day, the 
    maximum snow line elevation $H_s^{max}$ (above which every pixel is snow covered), and the minimum snow line elevation $H_s^{min}$ 
    (below which every pixel is snow free) are computed. A missing pixel located above $H_s^{max}$ is thus reclassified as snow 
    covered, while one located below $H_s^{min}$ is labelled as snow free.
    \item \textit{Direct neighbour spatial filling}, during which each missing pixel is compared with its four adjacent pixels. If 
    at least three share the same class, that class is assigned to the missing pixel.
    \item \textit{Elevation weighted spatial filling}, in which the $3\times3$ neighbourhood surrounding each missing grid point is 
    examined. If any neighbouring pixel is snow covered and lies at a lower elevation than the central one, the missing pixel is 
    reclassified as snow covered.
    \item \textit{Pixel-wise temporal filling}, which consists of examining the subsequent values of a missing grid point until a 
    snow covered or snow free classification is found. This value is then assigned to the missing pixel.
\end{itemize}

After the sixth step, the dataset is entirely free of missing values, providing snow cover information on a daily basis across the 
whole Italian peninsula.



\subsection{Results}

Following the approach adopted by the authors, particular attention was devoted to those metrics that describe the timing and the duration of the snow 
season. In particular, the snow cover duration (SCD) for each pixel was computed as
\begin{equation}
    \text{SCD}_{i,j}\,=\,\sum_{k=1}^{n} \text{S}_{i,j,k},
    \label{eq:scd}
\end{equation}
thereby summing, over a hydrological year spanning from the $1^{st}$ of October to the $30^{th}$ of September, the binary values recorded at each 
specific location. In Equation \eqref{eq:scd}, $k$ denotes the day of the hydrological year, whereas $(i,\,j)$ represent the spatial coordinates of the 
pixel. Two further snow metrics considered in the analysis of this dataset are the snow onset date (SOD) and the snow end date (SED). The former is 
defined as 
\begin{equation}
    \text{SOD}_{i,j}\,=\,\text{D}\,-\,\sum_{k=1}^{D} \text{S}_{i,j,k},
    \label{eq:sod}
\end{equation}
where $\text{D}$ corresponds to 116, namely the number of  days between the beginning of the hydrological year and the $25^{th}$ of 
January. Conversely, the latter one (SED) can be computed pixel-wise as
\begin{equation}
    \text{SED}_{i,j}\,=\,\text{D}\,+\,\sum_{k=D}^{n} \text{S}_{i,j,k},
    \label{eq:sed}
\end{equation}
with D retaining the same meaning as in Equation \eqref{eq:sod}. In what follows, only the results for SCD are presented: for the analogous analysis of 
SOD and SED, the reader is referred to Appendix \ref{app:MODIS_app}.

The mean snow cover duration across Italy, computed as an average over every pixel and over the 2001-2025 period, is 22 days, with an interannual standard deviation 
of 5 days. This rather low value is significantly biased by the abundance of plains and low-elevation terrains, which in the climatic setting of the study 
region remain essentially snow free all year round. MODIS observations indeed reveal that the snowpack is characterized by pronounced spatial variability: 
the longest snow seasons occur in high-altitude mountainous catchments, whereas the lowlands receive only sporadic and ephemeral snowfall. In the Po Plain, 
for instance, elevations below $\SI{500}{\meter}$ a.s.l. experience less than \num{5} days of snow cover per year on average, while the surrounding peaks 
are often snow covered for over \num{250} days a year. The mean SCD map, shown in Figure \ref{fig:ave_scd_MODIS}, closely follows the orography of 
the main mountain ranges, with valleys exhibiting markedly lower snow cover durations than the adjacent slopes. 
\begin{figure}[htb]
    \centering
    \includegraphics[width=\linewidth]{snow_products/ave_los_MODIS.png}
    \caption{Average SCD map across the Italian peninsula for the 2001-2025 period. SCD values accurately follow the orography of the 
    Italian Alps, as altitude strongly influence the persistence of snow on the ground. The lowest SCD values are associated with 
    lowland plains, which remain snow free for most of the year.}
    \label{fig:ave_scd_MODIS}
\end{figure}
It also underlines that the snow season in the Appenines is substantially shorter than in the Alps, with only the most prominent massifs, such as Gran 
Sasso and Maiella, displaying snow behaviour somewhat comparable to their Alpine counterpart. SOD and SED exhibit a similar spatial variability, as 
mountainous terrain strongly increases the likelihood of an earlier onset and a later end of the snow season. 

The predominant role of elevation in determining the duration of snow cover can be investigated by dividing the Italian territory into $\SI{5}{\meter}$ 
wide elevation bands. The quantiles of the SCD distribution, shown in Figure \ref{fig:scd_percentiles_by_ele_MODIS}, exhibit a clear monotonic 
relationship with elevation, as higher regions are typically associated with longer-lasting snowpacks. However, the considerable spread observed within 
each elevation band suggest that additional factors, such as climatic regimes or additional topographic variables, may influence snow cover duration across 
the Italian peninsula. For instance, repeating the same analysis for individual sub-regions of the investigated domain shows that the steepness of the 
quantile curve depends on the specific region considered. 
\begin{figure}[htpb]
    \centering
    \includegraphics[width=\linewidth]{snow_products/scd_percentiles_by_ele_MODIS.png}
    \caption{Quantiles of the average SCD distribution in each $\SI{5}{\meter}$ wide elevation band for the 2001-2025 period. The area 
    taken into account is the whole Italian peninsula. A clear monotonic trend with elevation is visible.}
    \label{fig:scd_percentiles_by_ele_MODIS}
\end{figure}
The snow regimes in the Alps and the Eastern Apennines are similar, with a rapid increase of 
SCD with elevation, whereas the Western Apennines, Sicily and Sardinia show a more gradual trend. Moreover, a multilinear regression with respect to 
several topographic and geographic parameter shows that aspect and latitude also provide a non-negligible contribution to snowpack properties \parencite{MODIS_ITA}.



The importance of aspect for the duration of the snow season is also evident when examining the spatial distribution of the 
difference between the SCD value of each individual pixel and the mean SCD per elevation band. In fact, in Figure 
\ref{fig:los_difference_by_ele_MODIS} the Alpine valleys are clearly delineated, since  south facing slopes experience, on average, 
a shorter snow season than north facing ones at the same elevation. The Easter Apennines can also be seen to be markedly snowier 
than the western side of the same ridge. This is attributed to cold continental air masses originating from the Balkan region, which 
become enriched with moisture as they cross the Adriatic Sea. This water vapor is subsequently released as snowfall through 
orographic forcing \parencite{more_snow_east_apennines}. This east to west gradient progressively weakens toward the Southern 
Apennines, which typically exhibit a shorter snow cover duration than the other Italian mountain ranges.

\begin{figure}[H]
    \centering
    \includegraphics[width=\linewidth]{snow_products/los_difference_by_ele_MODIS.png}
    \caption{Spatial distribution of the difference between the SCD value of each individual pixel and the mean SCD per $\SI{5}{\meter}$ 
    wide elevation band.}
    \label{fig:los_difference_by_ele_MODIS}
\end{figure}

The snow cover extent and duration observed in any given year cannot, of course, be expected to closely match the 2001-2025 
average, since the snowpack is highly sensitive to climatic drivers, such as winter precipitation and temperature patterns, that are 
themselves subject to considerable interannual variability. This translates into substantial year to year fluctuations in snow cover 
duration, which can be quantified through the pixel by pixel standard deviation from the mean, toghether with its ratio to the 
corresponding average SCD value. The spatial distribution of these two quantities is shown in Figure \ref{fig:std_over_mean_MODIS}. 
The standard deviation exhibits a monotonically increasing relationship with up to approximately $\SI{2000}{\meter}$ a.s.l., beyond 
which it begins to decrease. This behaviour may be partially linked to glacier covered terrain, which, as previously noted, is 
treated by MODIS as snow covered throughout the whole year, even though bare ice conditions are commonly observed by the end of the 
melt season. The ratio displays a markedly different pattern, with lower values typically found in high alpine catchments, whereas 
the highest ones are generally observerd over lowland plains. This trend follows directly from the considerably longer SCD 
characterising mountainous terrain, which contraints the ratio regardless of the absolute variability. 

% Plot va sistemato (unità di misura)!
\begin{figure}[H]
    \centering
    \includegraphics[width=\linewidth]{snow_products/std_over_mean_MODIS.png}
    \caption{(a) Spatial distribution of SCD standard deviation from the mean, computed over the 2001-2025 period. (b) Spatial 
    distribution of the ratio between the SCD standard deviation from the mean and the corresponding mean SCD value for the 2001-2025 
    period.}
    \label{fig:std_over_mean_MODIS}
\end{figure}










\section{IT-Snow}
\label{sec:it_snow}

IT-SNOW is a high resolution, multi-year snow reanalysis covering the entire Italian peninsula at a spatial resolution of roughly $\SI{500}{\meter}$. 
This dataset consist of daily maps of snow water equivalent (SWE), snow depth (HS), bulk snow density and liquid water content of the snowpack, spanning 
the period from \num{1} September \num{2010} to \num{31} August \num{2025} in its latest release. IT-SNOW is the output of S3M Italy, a real-time 
operational monitoring chain developed by CIMA Research Foundation and specifically designed for large scale applications, with the aim of serving as a 
valuable tool within flood forecasting frameworks. This model achieves a high level of performance by focusing on the key drivers of snow evolution, 
namely total precipitation, air temperature, relative humidity, and total radiation, thereby providing a spatially consistent reconstruction of snowpack 
dynamics.



\subsection{S3M}
\label{subsec:s3m}

S3M, also known as \textit{Snow Multidata Mapping and Modelling}, is a hydrology oriented cryospheric model developed to simulate the spatial and temporal 
evolution of seasonal snowpacks. It adopts a raster based approach, in which neighbouring pixels are treated as independent units, so that horizontal mass 
and energy transport is neglected. All input and output variables are therefore spatially distributed, and the simulated snowpack represents the average 
conditions within each grid cell. The model provides daily (or even hourly if needed) rasters of snow water equivalent (SWE), snow depth (HS), bulk snow 
density and liquid water content. These outputs can then be used to compute, during post-processing, snow cover maps spanning across the investigated 
domain, which could be used for direct comparison with remote sensing products. The snow cover fraction of each pixel can indeed be evaluated as
\begin{equation}
    \text{SCF}\,=\,\text{min}\left(1,\,\frac{\text{SWE [mm]}}{0.1 \cdot \rho_{s}}\right),
    \label{eq:scf_ITSNOW}
\end{equation}
where $\rho_{s}$ is the bulk snow density \parencite{scf_ITSNOW}. Since a pixel is conventionally regarded as snow covered when at least half of its area 
is, binary snow cover maps can be obtained by assigning a value of one to all grid points with $\text{SCF}>0.5$, and zero elsewhere.

Within S3M, seasonal snow is partitioned into a dry component, corresponding to the ice matrix, and a wet component, given by the interstitial liquid 
water. These two constituents are modelled separately, and their mutual interactions govern the overall snowpack dynamics. The total snow water equivalent 
is thus obtained as 
\begin{equation}
    \text{SWE}\,=\,\text{SWE}_D\,+\,\text{SWE}_W, 
    \label{eq:swe_tot_S3M}
\end{equation}
so as the sum of its dry ($\text{SWE}_D$) and wet ($\text{SWE}_W$) contributions. The evolution of these two quantities is described by the following mass 
conservation equations:
\begin{equation}
    \frac{d}{dt}\left(\text{SWE}_D\right)\,=\,S_f\,-\,M\,+\,R,
    \label{eq:swe_dry_S3M}
\end{equation}
\begin{equation}
    \frac{d}{dt}\left(\text{SWE}_W\right)\,=\,R_f\,+\,M\,-\,R\,-\,O,
    \label{eq:swe_wet_S3M}
\end{equation}
where $S_f$ and $R_f$ denote the snowfall and rainfall mass fluxes, $M$ and $R$ the snowmelt and refreezing rates and $O$ the snowpack runoff mass flux. 
All the terms in Equations \eqref{eq:swe_dry_S3M} and \eqref{eq:swe_wet_S3M} are expressed in $\SI{}{\milli\meter}$ w.e. $\SI{}{\per\hour}$ \parencite{S3M_model}. 
Solving this system of ordinary differential equations (ODEs) is a key step in simulating snowpack dynamics: in order to do so it is essential 
to make specific parametrizations for each of the mass fluxes involved.

Snowfall and rainfall mass fluxes are estimated from the total precipitation $P$ through a phase partitioning scheme based on air 
temperature and relative humidity. These two quantities are respectively computed as 
%\begin{equation}
%        S_f\,=\,\left(1\,-\,p_r\right) P
%    \label{eq:snowfall_S3M}
%\end{equation}
%\begin{equation}
%        R_f\,=\,p_rP
%    \label{eq:rainfall_S3M}
%\end{equation}
\begin{equation}
    \begin{cases}
        S_f\,=\,\left(1\,-\,p_r\right) P \\
        R_f\,=\,p_rP
    \end{cases},
    \label{eq:snowfall_rainfall_S3M}
\end{equation}
where $p_r$ is the rainfall probability, defined as
\begin{equation}
    p_r\,=\,\frac{1}{1\,+\,e^{\alpha\,+\,\beta T_{air}\,+\,\gamma RH}}.
    \label{eq:rain_probability_S3M}
\end{equation}
In the previous equation, $ T_{air}$ is the air temperature in $\SI{}{\degreeCelsius}$ and $RH$ is the relative humidity in percent, 
while the fixed parameter $\alpha\,=\,22$, $\beta\,=\,-2.7$ and $\gamma\,=\,-0.2$ are taken from \cite{S3M_phase_partitioning}. 

Snowmelt occurs only when both the current air temperature and its 10-day running mean exceed a threshold $T_{\tau}$, which S3M is set by default to 
$\SI{1}{\degreeCelsius}$. This dual condition provides a simple way of accounting for the snowpack's cold content, since by requiring sustained 
mild conditions, rather than a single warm hour, the model avoids triggering unphysical melt events that the snowpack's thermal inertia would not allow. 
S3M estimates the melt rate through a hybrid temperature index and radiation driven melt parametrization. This formulation improves the model's ability 
to reproduce the strong spatial variability induced by topography on snow cover extent and duration, since slope and aspect play a significant role in 
mountainous terrain that cannot be fully captured by air temperature alone. The hybrid approach adopted in S3M thus reads 
\begin{equation}
    M\,=\,m_{rad}\left[\frac{\left(1\,-\,\alpha\right)\text{SW}\downarrow}{\rho_W \lambda_f}\right]\Delta t\,+\,c_M m_r \left(T_{air}\,-\,T_{\tau}\right),
    \label{eq:melt_S3M}
\end{equation}
where $m_{rad}$ is a dimensionless coefficient that convert the potential radiative melt into actual melt, $\alpha$ is the snow 
albedo, $\lambda_f$ the specific latent heat of fusion and $m_r$ the degree day parameter [$\SI{}{\milli\meter\per\degreeCelsius\per\hour}$]. 
Both $m_{rad}$ and $m_r$ are calibration based.

In S3M, refreezing occurs whenever the air temperature drops below $T_{\tau}$ and is computed through a degree day approach as 
\begin{equation}
    R = - c_M m_r \left(T_{air}\,-\,T_{\tau}\right).
    \label{eq:refreezing_S3M}
\end{equation}
Since this expression depends solely on air temperature, refreezing is assumed to respond instantaneously to atmospheric 
cooling, without accounting for the internal structure or thermal gradient of the snowpack.

Ultimately, the runoff $O$ is estimated by simulating liquid water movement through the snowpack using a matrix-flow approximation. 
This liquid flux is governed by physical properties of the snow cover, as 
\begin{equation}
    O = \alpha \frac{\rho_W g K_W}{\mu_W}, 
    \label{eq:runoff_S3M}
\end{equation}
where $\alpha$ is a conversion parameter, $g$ the acceleration due to gravity, $K_W$ the intrinsic permeability of water in snow and 
$\mu_W$ is water dynamic viscosity. This parametrization completes the system of ODEs that, together with those governing the 
evolution of snow density, are solved at each grid point using the forward Euler method. By default, S3M adopts an hourly 
integration time step, although this can be adjusted by the user to meet specific application needs.


\subsection{Results}

According to the IT-SNOW reanalysis over the 2011-2025 period, the total volume of water stored within the snowpack on the national scale is on average 
$\SI{11.9 \pm 5.2}{G\meter\cubed}$, with a maximum of $\SI{23.1}{G\meter\cubed}$ in 2014 and a minimum of $\SI{5.4}{G\meter\cubed}$ in 2022. This peak is 
typically reached on 4 March $\pm$ 17 days, with the earliest occurrence being on 26 January 2023 and the latest on 26 March 2013. This means that over 
the investigated period the peak SWE condition is always achieved before \num{1} April, which is the date commonly adopted in literature as a reference 
for maximum SWE. The temporal evolution of the total SWE, which is shown in Figure \ref{fig:evo_swe_tot_ITSNOW}, exhibits a great year-to-year variability.
\begin{figure}[htpb]
    \centering
    \includegraphics[width=\linewidth]{snow_products/evo_swe_tot_ITSNOW.png}
    \caption{Daily evolution of the total water volume stored in the Italian snowpack, computed as snow water equivalent (SWE) 
    integrated over the whole national domain, from 1 September 2010 to 31 August 2025. Data are from the IT-SNOW reanalysis 
    \parencite{IT-SNOW}.}
    \label{fig:evo_swe_tot_ITSNOW}
\end{figure}
In particular, the last years of the record stand out for SWE values markedly below the long term mean, with the partial exception of 2024 since a 
unusually wet spring helped to balance the deficits accumulated during the early snow season. This pattern provides further evidence of the persistent 
snow drought conditions which have been detected in the Italian Alps during the early 2020s, driven by the combined effects of reduced precipitation and 
higher temperatures \parencite{snow_drought}. 
Besides a strong temporal variability, the modelled snow cover also displays a substantial spatial heterogeneity across Italy, as shown in Figure 
\ref{fig:ave_swe_scd_ITSNOW.png}. 
\begin{figure}[H]
    \centering
    \includegraphics[width=\linewidth]{snow_products/ave_swe_scd_ITSNOW.png}
    \caption{(a) Average seasonal snow water equivalent (SWE), computed as the mean of the November--June averages of each 
    hydrological year over the period 1 September 2010--31 August 2025. (b) Average snow cover duration (SCD), computed as the mean 
    of the annual SCD maps across the same period.}
    \label{fig:ave_swe_scd_ITSNOW.png}
\end{figure}
The simulated snowpack results indeed deeper and longer-lasting in high altitude catchments, whereas the lowlands only experience an ephemeral snow cover, 
owing to their milder temperatures. Water storage is also unevenly distributed between mountain ranges, since according to IT-SNOW estimates almost 87 \% 
of the national SWE is found within the five main Alpine river basins, namely the Po, Adige, Piave, Tagliamento and Brenta. The Central Apennines are 
Italy's second most important water tower, accounting for about 7\% of total SWE, whereas the remaining fraction is distributed across the rest of the 
country \parencite{IT-SNOW}. 













\section{Po River District}
\label{sec:po_district}

This snow dataset focuses on the mountainous portion of the Po River District, an over-regional entity that cover almost $\SI{87000}{\kilo\meter\squared}$ and comprises 
the Po River Basin together with some smaller adjacent catchments, such as Conca and Reno. This snow product is generated by an operational chain known 
as "J-snow", which is a hybrid approach that combines a physically based snow model named GEOtop \num{2.0} with in situ snow depth measurements and 
remote sensing snow cover extent estimates. SWE maps with a spatial resolution of approximately $\SI{500}{\meter}$ are available on a daily basis for all the hydrological 
years between 1992 and 2021. Within each hydrological year, maps are available only between October $3^{rd}$ and July $1^{st}$, a period that roughly 
corresponds to the snow season in non glacierized areas.

\subsection{GEOtop 2.0}

GEOtop 2.0 is a physically based model that simulates the heat and water exchanges at and below the ground surface. Its primary focus is thus the 
accurate reproduction of soil dynamics, including freezing and thawing processes as well as snow accumulation and melting, whenever thermal and 
meteorological conditions allow. 

The simulated soil is discretized into several cells, which are typically thinner near the ground surface in order to model as accurately as possible the 
strong interactions with the atmosphere taking place in the upper layers. The depth of the investigated domain is a user-defined parameter, which could 
range from a few meters to several hundred meters. The soil is treated with a rigid-soil scheme, meaning that expansion of the ground matrix is not 
allowed because GEOtop 2.0 lacks a mechanical model of porous media. This numerical compromise implies that no volume change is allowed 
during freezing and thawing processes and, consequently, the densities of liquid and solid water are both set at $\SI{1000}{\kilo\gram\per\meter\cubed}$. 
The simulated snowpack is directly coupled to the modelled soil, since the snow cover buffers the exchanges of energy and mass between the atmosphere and 
the ground. The heat and water content equations of the two media are solved as a single system by means of a time-lagged approach, whereby the coupling 
terms at time $t$ are replaced by their values at the previous time step $t-1$. This numerical method reduces the complexity of the problem, as each 
equation can now be solved separately, while preserving stability and first order accuracy. At each timestep, the snowpack is simulated by computing in 
sequence:

\begin{itemize}[label=$\blacktriangleright$]
    \item \textit{heat exchanges}, whose treatment follows the general framework presented in Chapter \ref{cap:snow_basics}, since in GEOtop 2.0 internal 
    energy varies under the forcing of shortwave and longwave radiative budgets, turbulent heat fluxes, precipitation heat fluxes and ground heat 
    fluxes. A distinctive feature of this model is that these terms are computed taking into account the complexity of the terrain, which is particularly 
    relevant in mountainous areas. The direct component of the shortwave radiation is indeed corrected for shading, since the surrounding topography can 
    hide the simulated pixel from the Sun, whereas the diffuse SW component is computed as the sum of the solar radiation scattered by the atmosphere and 
    of the shortwave radiation reflected by the portion of terrain visible from the pixel. Analogous corrections are applied to the thermal radiation, 
    which in mountainous regions typically originates both from the atmosphere and from the surrounding slopes. The topography also affects turbulent heat 
    fluxes, since it induces localized wind patterns resulting in a marked variability of latent and sensible heat exchanges. 
    \item \textit{snow metamorphism}, restricted to its destructive form. Constructive metamorphism, which is described in detail in Chapter \ref{cap:snow_basics}, 
    is not taken into account, as it is mainly relevant to avalanche forecasting and danger assessment, a purpose that differs considerably from the 
    hydrological one of this model. Compaction can therefore only result from destructive metamorphism acting on newly fallen snow and from the overburden 
    exerted by the overlying snowpack, so that the thickness of a snow layer evolves according to 
    \begin{equation}
        \frac{1}{D}\frac{\partial D}{\partial t}\,=\,C_1\,+\,C_2,
        \label{eq:layer_thickness_change_GEOtop}
    \end{equation}
        where $D$ is the vertical extent of each grid cell, $C_1$ is the compaction rate resulting from destructive metamorphism and $C_2$ is the compaction 
    rate caused by the overlying snow. Equation \eqref{eq:layer_thickness_change_GEOtop} is integrated with the same time step of the heat equation.
    \item \textit{water percolation}, which can be evaluated by solving a mass conservation equation. The water balance of each snow layer is indeed given 
    by 
    \begin{equation}
    \frac{\partial \theta_{\mathrm{w}}^{\mathrm{fl}}}{\partial t}
    + \nabla \cdot \boldsymbol{J}_{\mathrm{w}} + S_{\mathrm{w}} = 0 ,
    \label{eq:liquid_water_content}
    \end{equation}
    where $\theta_{\mathrm{w}}^{\mathrm{fl}}$ is the liquid water content, $\boldsymbol{J}_{\mathrm{w}}$ the flux of liquid water and $S_{\mathrm{w}}$ the 
    liquid water sink term. No percolation occurs until the liquid water content exceeds the saturation threshold of $3\%$ to $5\%$ of the snow layer mass
    (see Chapter \ref{cap:snow_basics}). Refreezing in sub-freezing layers is also accounted for: the internal energy carried by the water flux is tracked, 
    and latent heat is released whenever refreezing occurs.
    \item \textit{snow accumulation}, which is computed by means of a two air temperature threshold approach. Precipitation is assumed to fall entirely as 
    rain above the upper threshold and entirely as snow below the lower one, while in between the liquid and solid phases coexist and the solid fraction is 
    obtained by linear interpolation. The new snow density depends on both air temperature and wind speed, since it is evaluated as 
    \begin{equation}
        \rho_{\mathrm{ns}} =
        \begin{cases}
            500 - 475.5\exp\!\left[-1.4(5.0 - T)^{-1.15}
            - 8 \cdot \dfrac{v_{\mathrm{w}}^{1.7}}{1000}\right] \\[1ex]
            \qquad \text{if } T > -13^{\circ}\mathrm{C} \\[3ex]
            500 - 452.0\,\exp\!\left[-8 \cdot \dfrac{v_{\mathrm{w}}^{1.7}}{1000}\right] \\[1ex]
            \qquad  \text{elsewhere,}
        \end{cases}
        \label{eq:rho_swe_snow_GEOtop}
    \end{equation}
    where $T$ is the air temperature and $v_w$ is the wind speed.
\end{itemize}



\subsection{Results}

The total volume of water stored within the snowpack shows substantial year-to-year variability, as shown in Figure \ref{fig:evo_swe_tot_PO}, with 
an average peak value of $6.7 \pm 2.3 \,\SI{}{\giga\meter\cubed}$ typically reached on $20 \text{March} \,\pm\,25$ days. This date is compatible within 
one standard deviation with the 1\textsuperscript{st} of April, which is the date conventionally used as a proxy for maximum SWE in snow literature.

\begin{figure}[htbp]
    \centering
    \includegraphics[width=\linewidth]{snow_products/evo_swe_tot_PO.png}
    \caption{Daily evolution of the total water volume stored within the snowpack in the mountainous portion of the Po River District. The investigated 
    period spans from 3 October 1991 to 1 July 2021. Data are from the Po River District snow dataset \parencite{Po_snow}}
    \label{fig:evo_swe_tot_PO}
\end{figure}

The vast majority of this peak SWE is stored above

\begin{table}[tb]
  \centering
  \caption{Po Basin}
  \label{tab:peak_swe_elevation_wise_PO}
  \renewcommand{\arraystretch}{1.3}
  \renewcommand{\tabularxcolumn}[1]{m{#1}}   % X centrate verticalmente
  \setlength{\tabcolsep}{6pt}
  \footnotesize
  \begin{tabularx}{\textwidth}{@{}
      >{\raggedright\arraybackslash}m{3.0cm}
      >{\raggedright\arraybackslash}m{3.0cm}
      >{\raggedright\arraybackslash}m{3.0cm}
      >{\raggedright\arraybackslash}X @{}}
    \toprule
    \textbf{Elevation band} & \textbf{Peak SWE} & \textbf{Peak date} & \textbf{SWE trend} \\
    \midrule
    0 - 500 m & $0.08 \pm 0.06 \,\SI{}{\giga\meter\cubed}$ & 10 Jan $\pm$ 33 days &  $-\,\SI{0.01}{\giga\meter\cubed}$/decade (p = 0.10)\\
    500 - 1000 m & $0.22 \pm 0.11 \,\SI{}{\giga\meter\cubed}$ & 13 Jan $\pm$ 32 days & $-\,\SI{0.02}{\giga\meter\cubed}$/decade (p = 0.41)\\
    1000 - 1500 m & $0.63 \pm 0.39 \,\SI{}{\giga\meter\cubed}$ & 04 Feb $\pm$ 29 days & $-\,\SI{0.06}{\giga\meter\cubed}$/decade (p = 0.30)\\
    1500 - 2000 m & $1.6 \pm 0.8 \,\SI{}{\giga\meter\cubed}$ & 06 Mar $\pm$ 23 days & $+\,\SI{0.06}{\giga\meter\cubed}$/decade (p = 0.50)\\
    2000 - 2500 m & $2.7 \pm 0.9 \,\SI{}{\giga\meter\cubed}$ & 05 Apr $\pm$ 19 days & $+\,\SI{0.21}{\giga\meter\cubed}$/decade (p = 0.27)\\
    > 2500 m& $2.3 \pm 0.6 \,\SI{}{\giga\meter\cubed}$ & 30 Apr $\pm$ 17 days & $+\,\SI{0.19}{\giga\meter\cubed}$/decade (p = 0.20)\\
    \bottomrule
  \end{tabularx}
\end{table}






\section{Comparison}

The snow datasets presented above differ from one another in several respects. ITSNOW and the Po Basin product are the outputs of two distinct 
physically based model, and both focus primarily on snow water equivalent, consistently with their hydrological purpose. ITSNOW, however, provides 
richer information than the latter, as it also includes maps of snow density, snow depth and liquid water content, which allow the user to gain 
a deeper insight into the state of the snowpack. MODIS, in contrast, describes snow cover alone through binary rasters. This discrepancy does not 
limit the possibility of a comparison, as SWE maps can be converted into binary rasters by means of \eqref{eq:scf_ITSNOW} and \eqref{eq:scf_PO}. Since 
these datasets cover different domains, are provided at different spatial resolutions and in different coordinate reference systems, and span 
different periods, the comparisons reported hereafter are carried out over the largest common overlap, both in space and in time. For the sake of 
clarity, Table \ref{tab:dataset_car} summarizes the main characteristics of each snow product, thereby clarifying the rationale behind some of the 
comparison choices.

\begin{table}[H]
  \centering
  \caption{Main characteristics of each snow product.}
  \label{tab:dataset_car}
  \renewcommand{\arraystretch}{1.3}
  \renewcommand{\tabularxcolumn}[1]{m{#1}}   % X centrate verticalmente
  \setlength{\tabcolsep}{4pt}
  \footnotesize
  \begin{tabularx}{\textwidth}{@{} >{\raggedright\arraybackslash}m{1.9cm} >{\raggedright\arraybackslash}X >{\raggedright\arraybackslash}m{2.1cm} >{\raggedright\arraybackslash}m{2.1cm} >{\raggedright\arraybackslash}m{1.5cm} >{\raggedright\arraybackslash}m{2.1cm} >{\raggedright\arraybackslash}X @{}}
    \toprule
    \textbf{Dataset} & \textbf{Data type} & \textbf{Start date} & \textbf{End date} & \textbf{Study area} & \textbf{Spatial resolution} & \textbf{Reference system} \\
    \midrule
    MODIS & Snow cover & 24-02-2000 & 30-09-2025 & Italy & 0.00417 $\unit{\degree}$ &  WGS 84 (EPSG:4326) \\
    ITSNOW & SWE & 01-09-2010 & 31-08-2025 & Italy & 0.00506 $\unit{\degree}$ & WGS 84 (EPSG:4326) \\
    Po Basin & SWE & 03-10-1991 & 01-07-2021 & Po Basin & $\SI{500}{\meter}$ & ETRS89 (EPSG:3035) \\
    \bottomrule
  \end{tabularx}
\end{table}

\begin{figure}[H]
    \centering
    \includegraphics[width=\linewidth]{snow_products/comparison_SCD_MODIS_ITSNOW.png}
    \caption{MODIS}
    \label{fig:comparison_SCD_MODIS_ITSNOW}
\end{figure}
